\documentclass[11pt]{article}
\usepackage[margin=2.2cm]{geometry}
\usepackage{fontspec}
\setmainfont{cmun}[Extension=.otf, UprightFont=*rm, BoldFont=*bx, ItalicFont=*ti, BoldItalicFont=*bi]
\setsansfont{cmun}[Extension=.otf, UprightFont=*ss, BoldFont=*sx, ItalicFont=*si, BoldItalicFont=*so]
\setmonofont{cmun}[Extension=.otf, UprightFont=*tt, BoldFont=*tb, ItalicFont=*it, BoldItalicFont=*tx]
\usepackage[russian]{babel}
\usepackage{amsmath,graphicx,booktabs,xcolor,hyperref,enumitem,caption}
\hypersetup{colorlinks=true,linkcolor=blue!50!black,citecolor=blue!50!black,urlcolor=blue!50!black}
\setlist{itemsep=2pt,topsep=4pt}
\newcommand{\done}{\textcolor{green!50!black}{\textbf{готово}}}
\newcommand{\running}{\textcolor{orange!80!black}{\textbf{идёт}}}
\newcommand{\todo}{\textcolor{gray}{\textbf{впереди}}}
\newcommand{\OverviewResults}{\textcolor{gray}{(результаты появятся после завершения сетки: 3 семейства, 3 сида, контроли и масштабирование)}}


\title{\textbf{VIBES: world model учит физику или запоминает картинки?}\\[0.3em]
\large Идея исследования, что нашли, куда подавать}
\author{nonberw}
\date{ноябрь 2026 г.}

\begin{document}
\maketitle

\begin{center}
\fbox{\parbox{0.92\linewidth}{\textbf{Коротко.} Видео-модели мира (world models) предсказывают будущее,
которое выглядит правильно, пока физика такая же, как при обучении. Стоит изменить, например,
гравитацию за пределы обучающего диапазона, и они ломаются. Все это видят по выходу модели. Мы
спрашиваем, \emph{где именно} внутри модели это ломается: в \textbf{выводе} (модель не может
поверить в непривычную гравитацию), в \textbf{исполнении} (динамика не умеет такую гравитацию
разыграть) или в \textbf{восприятии}. И доказываем, что <<ломаться>> тут значит делать ровно то,
чего требует функция потерь.}}
\end{center}

\section{Проблема и что уже известно}

В 2024--2026 годах вышло много работ, которые проверяют физику видео-моделей по их выходу.
Kang et al.\ (ICML 2025) на 2D-симуляторе показали: внутри распределения почти идеально, вне
распределения провал, причём модель копирует ближайший обучающий пример. Physics-IQ и Morpheus
показывают то же на реальных видео. Другая линия работ (Joseph et al., Tan et al., Liao \& Cao,
Bao) смотрит внутрь моделей линейными пробами: физические величины там \emph{читаются}, но
часто \emph{не используются}.

Чего не хватает:
\begin{itemize}
\item поведение говорит, \emph{что} модель сломалась вне диапазона, но не \emph{где};
\item пробы говорят, что информация есть, но не что модель ей пользуется;
\item никто не спросил, что в этой ситуации сделал бы \emph{идеальный} обучаемый.
\end{itemize}

\section{Два теоретических факта (Раздел 3 статьи)}

\begin{enumerate}
\item \textbf{Компетентность влечёт декодируемость.} Если модель хорошо предсказывает будущее, то
гравитация и энергия обязаны быть восстановимы из её состояния. Поэтому положительная проба~---
слабое свидетельство <<понимания физики>>: она следует из низкой ошибки.
\item \textbf{Bayes-оптимальный предиктор клампится.} При квадратичной ошибке оптимальное
предсказание несёт в себе \emph{апостериорное среднее} гравитации, а оно всегда лежит внутри
обучающего диапазона. С шумом наблюдений оно сходится к границе диапазона. Без шума функция потерь
вообще ничего не говорит о контекстах вне диапазона. Вывод: <<прилипание к границе>>~--- это не
баг, а то, чего просит лосс; экстраполяция закона может взяться только из индуктивного смещения
архитектуры.
\end{enumerate}

\section{Как мы это измеряем}

\textbf{Мир.} Коробка с шарами и блоками, 1--3 объекта, свой $g\in[4,8]$, упругость $e$ и трение
$\mu$ в каждом эпизоде. Материалы \emph{видны} (оттенок фона), гравитация~--- нет, её надо выводить
из движения. 20 кадров/с, $64{\times}64$, 10 кадров контекста, 20 кадров предсказания, 100\,000
обучающих эпизодов. Тестовый набор \texttt{gsweep}: всё как при обучении, но $g$ от 1 до 12.

\textbf{Три семейства, одна динамика.} Кодировщик кадра обучается и замораживается, потом на его
латентах обучается одна и та же chunk-модель динамики (belief $h_C$ по контексту, декодер сразу на
20 шагов вперёд). Семейства отличаются только сигналом обучения кодировщика:
REC (реконструкция пикселей), JEPA (предсказание латентов, без декодера), STATE (истинные
координаты, без скоростей и без $g$).

\textbf{Три способа померить гравитацию в модели.}
\begin{enumerate}
\item \emph{Естественный отклик}: какую гравитацию показывают предсказания модели.
\item \emph{Внутренняя оценка}: что про $g$ читает линейная проба из belief'а.
\item \emph{Запись}: сдвигаем belief вдоль направления гравитации и смотрим, что исполнит
динамика. Направление берём из \emph{точных контрфактических пар}: обращаем симулятор во времени и
строим два контекста, которые заканчиваются в одном и том же состоянии, но при разной гравитации.
\end{enumerate}

\textbf{Вложенная инверсия (главный инструмент).} Линейный сдвиг может не сработать по скучной
причине (гравитация закодирована криво). Поэтому мы \emph{оптимизируем} вход замороженной модели:
сам belief (тест исполнения) или латенты контекста, пропущенные через замороженную belief-сеть
(тест достижимости). Подгоняем только первые 4 предсказанных кадра, а физику измеряем на
\emph{отложенных} кадрах 4--13. Контроль <<антифизики>> (гравитация вверх, $g=-4$) показывает,
насколько декодер вообще гибкий.

\textbf{Предрегистрация.} Гипотезы, метрики, правило решения и тесты закоммичены до оценок
(\texttt{PREREGISTRATION\_v2.md}); каждое отклонение записано с датой и причиной. Первая версия
исследования (\texttt{PREREGISTRATION.md}) провалилась: модели не выучили физику вообще. Мы это
честно описываем в приложении статьи, потому что причины провала (мало данных, скрытые материалы,
авторегрессия) сами по себе информативны.

\section{Результаты}
\label{sec:results}

\OverviewResults

\section{Чем это отличается от предыдущих работ}

\begin{itemize}
\item Kang et al.\ и Physics-IQ смотрят только на выход. Мы локализуем провал внутри модели
причинными вмешательствами.
\item Работы с пробами (Joseph, Tan, Liao \& Cao, Bao) показывают, что величины читаются или не
используются. Мы спрашиваем другое: может ли динамика вообще исполнить закон вне диапазона, и
достижим ли нужный belief.
\item Контрфактические пары с точно совпадающим состоянием (обращение времени в симуляторе) и
инверсия с отложенными кадрами и антифизическим контролем~--- новая методика.
\item Bayes-референс превращает <<провал>> в проверяемое предсказание теории.
\end{itemize}

\section{Куда подавать}

\begin{itemize}
\item \textbf{ICML 2027} (основной трек, Рио-де-Жанейро, 31 июля~--- 6 августа 2027)~---
рекомендуемая цель. Официальный call for papers опубликован: сайт подачи открывается 8 января 2027,
дедлайн абстракта \textbf{14 января 2027, 23:59 AoE}, полной статьи \textbf{22 января 2027, 23:59
AoE}, без продлений. 8 страниц основного текста плюс неограниченные ссылки и приложения, слепое
рецензирование, обязательны Impact Statement и AI Disclosure Statement (оба уже в статье). Статья
уже свёрстана официальным стилем \texttt{icml2027}.
\item \textbf{Воркшопы ICLR 2027} (Сан-Франциско, 29--30 апреля 2027)~--- запасной или
параллельный вариант. Список воркшопов объявят 29 ноября 2026, рекомендованный дедлайн статей
1 февраля 2027. Проверить правила двойной подачи.
\item CVPR 2027 (16 ноября 2026) и L4DC 2027 (13 ноября 2026) уже не успеть.
\end{itemize}

\section{Что нужно сделать автору}

Подача требует личных аккаунтов, мы ничего не отправляли. Пошаговый список:
\texttt{docs/SUBMISSION.md}. Главное: дождаться CFP ICML 2027, заменить стилевой файл, убрать имя и
ссылку на репозиторий для слепого рецензирования, добавить impact statement, загрузить на
OpenReview; при желании выложить на arXiv (cs.LG).

\end{document}
